\chapter{From Open-String Quanta to Tachyons and Gravity}

\lectureoverview{We now quantize the normal modes found in the previous
lecture and ask what particle states they produce. Each transverse mode is a
harmonic oscillator, but its normalization is slightly unfamiliar: exciting
mode \(n\) raises the light-front energy, and hence the mass-squared level, by
\(n\) units. The first excited open-string level transforms exactly like the
two transverse polarizations of a massless vector. Relativistic consistency
therefore fixes the ground-state mass squared to be negative. That state is
the bosonic-string tachyon, whose real significance is an unstable vacuum,
not a practical faster-than-light particle. Finally, once strings interact,
joining and splitting make closed strings unavoidable; the lowest closed-
string excitation will be the graviton.}

\section{The oscillator hidden in each string mode}

For one transverse coordinate \(X\), the \(n\)th nonzero normal mode has
energy and Lagrangian

\begin{align}
  E_n
  &=
  \frac14\dot X_n^{\,2}
  +\frac{n^2}{4}X_n^2,
  \label{eq:l03-mode-energy}\\
  L_n
  &=
  \frac14\dot X_n^{\,2}
  -\frac{n^2}{4}X_n^2.
  \label{eq:l03-mode-lagrangian}
\end{align}

The factors of \(1/4\) are not an error. They record the normalization chosen
when the string action was reduced to independent Fourier modes. The first
step in constructing the Hamiltonian is therefore to compute the canonical
momentum,

\begin{equation}
  p_n=\frac{\partial L_n}{\partial\dot X_n}
     =\frac12\dot X_n,
  \qquad
  \dot X_n=2p_n.
  \label{eq:l03-canonical-momentum}
\end{equation}

The Legendre transform gives

\begin{equation}
  H_n=p_n\dot X_n-L_n
     =p_n^2+\frac{n^2}{4}X_n^2.
  \label{eq:l03-mode-hamiltonian}
\end{equation}

\begin{figure}[H]
  \centering
  \includegraphics[width=0.82\textwidth]{lecture_03_figure_01.png}
  \caption{The mode energy, Lagrangian, canonical momentum, and Hamiltonian
  in the normalization inherited from the open string.}
  \lectureframe{00:05:40}
  \label{fig:l03-mode-hamiltonian-board}
\end{figure}

To factorize \(H_n\), combine \(X_n\) and \(p_n\) as

\begin{equation}
  \frac{nX_n}{2}\pm i p_n.
\end{equation}

Classically the two factors would commute, but quantum mechanically

\begin{equation}
  [X_n,p_m]=i\delta_{nm}
  \label{eq:l03-canonical-commutator}
\end{equation}

in units with \(\hbar=1\). The cross terms do not cancel innocently. For a
fixed mode,

\begin{align}
  \left[\frac{nX_n}{2}+ip_n,
        \frac{nX_n}{2}-ip_n\right]
  &=n,
  \label{eq:l03-unscaled-commutator}\\
  \left(\frac{nX_n}{2}-ip_n\right)
  \left(\frac{nX_n}{2}+ip_n\right)
  &=H_n-\frac n2.
  \label{eq:l03-factorized-hamiltonian}
\end{align}

\begin{figure}[H]
  \centering
  \includegraphics[width=0.82\textwidth]{lecture_03_figure_02.png}
  \caption{The board calculation in which the canonical commutator supplies
  the factor of \(n\).}
  \lectureframe{00:10:30}
  \label{fig:l03-commutator-board}
\end{figure}

The factor of \(n\) tells us how to normalize the ladder operators:

\begin{align}
  a_n
  &\equiv
  \frac{\sqrt n}{2}X_n+\frac{i}{\sqrt n}p_n,
  \label{eq:l03-annihilator}\\
  a_n^\dagger
  &\equiv
  \frac{\sqrt n}{2}X_n-\frac{i}{\sqrt n}p_n.
  \label{eq:l03-creator}
\end{align}

They obey

\begin{equation}
  [a_n,a_m^\dagger]=\delta_{nm},
  \qquad
  H_n=n\left(a_n^\dagger a_n+\frac12\right).
  \label{eq:l03-oscillator-algebra}
\end{equation}

Thus \(a_n^\dagger\) adds one quantum of the \(n\)th oscillator, but that
quantum carries \(n\) units of mode energy. The familiar oscillator spacing
is not lost; its frequency is \(\omega_n=n\).

\begin{figure}[H]
  \centering
  \includegraphics[width=0.82\textwidth]{lecture_03_figure_03.png}
  \caption{Normalized annihilation and creation operators for the \(n\)th
  string oscillator.}
  \lectureframe{00:12:00}
  \label{fig:l03-ladder-board}
\end{figure}

\subsection{Two transverse towers}

In four spacetime dimensions, a string carrying very large momentum along
the \(z\)-axis has two transverse coordinates, \(X\) and \(Y\). We therefore
need two independent oscillator towers. Let \(a_n,a_n^\dagger\) act on the
\(X\) modes and \(b_n,b_n^\dagger\) on the \(Y\) modes. Solving
\eqref{eq:l03-annihilator}--\eqref{eq:l03-creator} for the coordinates gives

\begin{equation}
  X_n=\frac{a_n+a_n^\dagger}{\sqrt n},
  \qquad
  Y_n=\frac{b_n+b_n^\dagger}{\sqrt n}.
  \label{eq:l03-mode-coordinates}
\end{equation}

\begin{figure}[H]
  \centering
  \includegraphics[width=0.78\textwidth]{lecture_03_figure_04.png}
  \caption{Solving for a mode coordinate in terms of its creation and
  annihilation operators.}
  \lectureframe{00:17:05}
  \label{fig:l03-coordinate-operator-board}
\end{figure}

The open-string coordinates were expanded in Neumann modes,

\begin{align}
  X(\sigma,\tau)
  &=X_{\rm cm}(\tau)
    +\sum_{n=1}^{\infty}X_n(\tau)\cos(n\sigma),
  \label{eq:l03-X-fourier}\\
  Y(\sigma,\tau)
  &=Y_{\rm cm}(\tau)
    +\sum_{n=1}^{\infty}Y_n(\tau)\cos(n\sigma).
  \label{eq:l03-Y-fourier}
\end{align}

Cosines appear because the endpoint slopes vanish. Substituting
\eqref{eq:l03-mode-coordinates} turns the string itself into an operator:

\begin{align}
  X(\sigma,\tau)
  &=X_{\rm cm}(\tau)
    +\sum_{n=1}^{\infty}
      \frac{a_n(\tau)+a_n^\dagger(\tau)}{\sqrt n}
      \cos(n\sigma),
  \label{eq:l03-X-operator}\\
  Y(\sigma,\tau)
  &=Y_{\rm cm}(\tau)
    +\sum_{n=1}^{\infty}
      \frac{b_n(\tau)+b_n^\dagger(\tau)}{\sqrt n}
      \cos(n\sigma).
  \label{eq:l03-Y-operator}
\end{align}

\begin{figure}[H]
  \centering
  \begin{minipage}[t]{0.48\textwidth}
    \centering
    \includegraphics[width=\textwidth]{lecture_03_figure_05.png}
    \par\smallskip\footnotesize (a) Coordinate-mode expansions.
  \end{minipage}\hfill
  \begin{minipage}[t]{0.48\textwidth}
    \centering
    \includegraphics[width=\textwidth]{lecture_03_figure_06.png}
    \par\smallskip\footnotesize (b) Ladder operators substituted into the
    string coordinates.
  \end{minipage}
  \caption{The shape of the quantum open string is assembled from all of its
  transverse oscillator operators.}
  \lectureframe{00:20:35 and 00:21:05}
  \label{fig:l03-string-coordinate-board}
\end{figure}

\begin{classroomqa}
\audiencequestion{Is \(X_n\) also a function of
\(\sigma\)?\lecturetimestamp{00:21:48}}
\lecturerresponse{No. The \(\sigma\)-dependence is entirely in
\(\cos(n\sigma)\). The coefficient \(X_n\) is the time-dependent amplitude of
that spatial pattern, so \(X_n=X_n(\tau)\).}
\end{classroomqa}

\begin{classroomqa}
\audiencequestion{Why do the operators called raising and lowering seem to
have their signs reversed?\lecturetimestamp{00:23:13}}
\lecturerresponse{The visible sign of \(ip_n\) is not a naming rule. Compute
the commutator with the Hamiltonian. The operator satisfying
\([H_n,a_n^\dagger]=n a_n^\dagger\) raises the energy; the one satisfying
\([H_n,a_n]=-n a_n\) lowers it.}
\end{classroomqa}

The two transverse towers are a pedagogical four-dimensional picture. A
bosonic string quantized consistently has twenty-six spacetime dimensions,
so it has twenty-four transverse oscillator families rather than two. The
logic below is unchanged: the number of first-level states must match the
number of physical polarizations of a massless vector in the same dimension.

\section{Massive spin and massless helicity}

Before assigning spin to the string states, we need one sharp distinction.
A massive particle of spin \(J\) has \(2J+1\) spin states,

\begin{equation}
  m_J=-J,-J+1,\ldots,J-1,J.
  \label{eq:l03-massive-spin-multiplet}
\end{equation}

For example, a massive spin-one particle has three polarizations and a
massive spin-two particle has five. We can see why by going to the rest
frame. There every spatial rotation is available. Rotate the spin axis,
boost back, and all \(2J+1\) components belong to the same relativistic
particle.

A massless particle has no rest frame. Its invariant spin label is helicity,
the angular momentum along its direction of motion. A photon has only
\(h=+1\) and \(h=-1\); a graviton has only \(h=+2\) and \(h=-2\). A scalar has
one \(h=0\) state. The lecture uses reflection symmetry to pair the two
helicities. In a chiral theory a single massless irreducible representation
may carry only one helicity, so the pairing is an additional symmetry
statement rather than a consequence of masslessness alone.\editorialnote{The
modern little-group classification permits one-helicity massless
representations. Parity or an appropriate CPT-complete field content supplies
the opposite helicity in the familiar photon and graviton examples.}

\begin{figure}[H]
  \centering
  \begin{minipage}[t]{0.45\textwidth}
    \centering
    \includegraphics[width=\textwidth]{lecture_03_figure_07.png}
    \par\smallskip\footnotesize (a) The spin-state count begun on the board.
  \end{minipage}\hfill
  \begin{minipage}[t]{0.5\textwidth}
    \centering
    \begin{tabular}{@{}ccc@{}}
      \toprule
      Spin & Massive states & Massless helicities \\
      \midrule
      \(0\) & \(1\) & \(0\) \\
      \(1\) & \(3\) & \(+1,-1\) \\
      \(2\) & \(5\) & \(+2,-2\) \\
      \bottomrule
    \end{tabular}
  \end{minipage}
  \caption{Mass removes the rest frame and with it the intermediate spin
  projections.}
  \lectureframe{00:27:40}
  \label{fig:l03-spin-count}
\end{figure}

\subsection{Linear and circular photon polarization}

Take a photon moving along \(z\). Its two transverse linear polarizations are
\(\ket{x}\) and \(\ket{y}\). They are not helicity eigenstates. The circular
states are

\begin{align}
  \ket{R}
  &=\frac{1}{\sqrt2}\left(\ket{x}+i\ket{y}\right),
  \label{eq:l03-right-circular}\\
  \ket{L}
  &=\frac{1}{\sqrt2}\left(\ket{x}-i\ket{y}\right),
  \label{eq:l03-left-circular}
\end{align}

with helicities \(+1\) and \(-1\), up to the convention used to name right
and left. A linearly polarized photon is a superposition of those two
outcomes. Its average helicity can vanish, but a helicity measurement never
finds a spin-zero photon. Likewise,

\begin{equation}
  \frac{1}{\sqrt2}\left(\ket{x}+\ket{y}\right)
\end{equation}

is simply linear polarization at forty-five degrees.

\begin{figure}[H]
  \centering
  \begin{minipage}[t]{0.48\textwidth}
    \centering
    \includegraphics[width=\textwidth]{lecture_03_figure_08.png}
    \par\smallskip\footnotesize (a) Circular polarization states on the
    board.
  \end{minipage}\hfill
  \begin{minipage}[t]{0.48\textwidth}
    \centering
    \begin{tikzpicture}[>=Latex,scale=0.9]
      \draw[->,thick] (0,0) -- (3.4,0) node[right] {\(z\)};
      \draw[->,blue!70!black,thick] (0,0) -- (0,2.0)
        node[above] {\(x\)};
      \draw[->,orange!85!black,thick] (0,0) -- (-1.4,-1.2)
        node[below left] {\(y\)};
      \draw[red!70!black,thick,->]
        plot[domain=0:320,samples=80]
        ({1.0*cos(\x)},{0.58*sin(\x)+0.4});
      \node[red!70!black] at (1.55,1.25) {circular};
    \end{tikzpicture}
  \end{minipage}
  \caption{Linear polarization supplies a transverse basis; circular
  combinations diagonalize helicity.}
  \lectureframe{00:37:45}
  \label{fig:l03-polarization}
\end{figure}

\begin{classroomqa}
\audiencequestion{Is this like the separation of spin states in a
Stern--Gerlach experiment?\lecturetimestamp{00:37:40}}
\lecturerresponse{The quantum-superposition idea is the same, but for the
photon the direct language is polarization. A circular-polarization
measurement distinguishes the two helicities; there is no longitudinal
spin-zero photon in the beam.}
\end{classroomqa}

The graviton repeats the pattern at spin two. For spin bookkeeping one may
write schematically

\begin{equation}
  \ket{+2}\sim\ket{+1}\otimes\ket{+1},
  \qquad
  \ket{-2}\sim\ket{-1}\otimes\ket{-1}.
  \label{eq:l03-graviton-helicity}
\end{equation}

This does not say that a graviton is literally made from two photons. It says
that two equal helicities add to the correct spin-two helicity.

\begin{classroomqa}
\audiencequestion{Does the two-photon analogy extend to other quantum
numbers?\lecturetimestamp{00:41:05}}
\lecturerresponse{Not as a constituent model. There are powerful formulas in
which gravity amplitudes look like products, or ``squares,'' of gauge-theory
amplitudes, but the graviton is still its own particle. The analogy is most
trustworthy here as a statement about helicity.\editorialnote{Modern
Kawai--Lewellen--Tye and double-copy relations make the amplitude statement
precise in suitable theories; they do not turn a graviton into a bound state
of two photons.}}
\end{classroomqa}

\begin{classroomqa}
\audiencequestion{A photon can have any momentum, so does it not have a
continuous spectrum rather than the isolated spectrum required of a
particle?\lecturetimestamp{00:42:17}}
\lecturerresponse{The continuous variable is the center-of-mass momentum.
Hold that momentum fixed and inspect the internal invariant mass. A particle
is identified by an isolated mass level; varying its overall momentum then
produces the usual continuous dispersion relation.}
\end{classroomqa}

\section{The lowest open-string levels}

We can now return to the string. In the large-momentum description, the
finite internal light-front energy is proportional to invariant mass
squared. Suppressing the overall string scale, write the spectrum as

\begin{equation}
  M^2
  =m_0^2
  +\sum_{n=1}^{\infty}n
   \left(N_n^{(a)}+N_n^{(b)}\right),
  \qquad
  N_n^{(a)}=a_n^\dagger a_n,
  \quad
  N_n^{(b)}=b_n^\dagger b_n.
  \label{eq:l03-open-spectrum}
\end{equation}

The oscillator ground state obeys

\begin{equation}
  a_n\ket{0}=0,
  \qquad
  b_n\ket{0}=0,
  \qquad n=1,2,\ldots .
  \label{eq:l03-ground-annihilation}
\end{equation}

It is not an empty absence of string. It is the least excited quantum string,
and its coordinates still fluctuate. Let its mass squared be \(m_0^2\), not
yet determined.

\begin{figure}[H]
  \centering
  \includegraphics[width=0.78\textwidth]{lecture_03_figure_09.png}
  \caption{The open-string ground state and the annihilation conditions that
  define it.}
  \lectureframe{00:45:20}
  \label{fig:l03-ground-state-board}
\end{figure}

The cheapest excitations use the \(n=1\) creation operators:

\begin{equation}
  a_1^\dagger\ket{0},
  \qquad
  b_1^\dagger\ket{0},
  \qquad
  M^2=m_0^2+1.
  \label{eq:l03-first-level}
\end{equation}

\begin{figure}[H]
  \centering
  \includegraphics[width=0.78\textwidth]{lecture_03_figure_10.png}
  \caption{One of the two degenerate first-level states, with its
  mass-squared increment written on the board.}
  \lectureframe{00:49:25}
  \label{fig:l03-first-level-board}
\end{figure}

The ``one'' in eqref{eq:l03-first-level} is not one joule, one electron volt,
or one Planck mass. We have suppressed the dimensionful string scale. For
now it means one step in the oscillator's mass-squared ladder; the missing
units will be restored in the next lecture.

Any linear combination of the two degenerate states is also a state at this
level. In particular,

\begin{equation}
  \frac{1}{\sqrt2}
  \left(a_1^\dagger\pm i b_1^\dagger\right)\ket{0}
  \label{eq:l03-string-circular-states}
\end{equation}

is the string analogue of circular polarization. Under a rotation in the
\(x\)-\(y\) plane, \(a_1^\dagger\) and \(b_1^\dagger\) transform as the two
components of a transverse vector. The combinations in
\eqref{eq:l03-string-circular-states} therefore have helicity \(+1\) and
\(-1\).

\begin{figure}[H]
  \centering
  \includegraphics[width=0.82\textwidth]{lecture_03_figure_11.png}
  \caption{Ground-state conditions, the two first-level excitations, and
  their circular combinations collected on the board.}
  \lectureframe{00:56:00}
  \label{fig:l03-level-structure-board}
\end{figure}

There is no third degenerate state that could serve as a longitudinal
polarization. If these two states had nonzero mass, a spin-one relativistic
particle would require three polarizations. With exactly two, the state must
be massless. Hence

\begin{equation}
  m_0^2+1=0,
  \qquad
  m_0^2=-1.
  \label{eq:l03-intercept-condition}
\end{equation}

The first level is photon-like: it is a massless gauge-vector state. It need
not be the photon of electromagnetism; its eventual gauge identity depends on
the string model. The ground state below it has negative mass squared. That
is the bosonic open-string tachyon.\editorialnote{In the full critical
bosonic string the normal-ordering intercept and Lorentz algebra are tied to
\(D=26\). The lecture's two-transverse-direction calculation is the
four-dimensional shadow of the same state-counting argument: \(D-2\)
first-level oscillator states match the \(D-2\) physical polarizations of a
massless vector.}

\begin{classroomqa}
\audiencequestion{Why should we care that \(m_0^2\) is
negative?\lecturetimestamp{00:56:35}}
\lecturerresponse{Because mass squared is the curvature of the quadratic
field potential about the state we have called the vacuum. Negative
curvature means that this configuration is not a stable ground state. The
next calculation makes that statement precise.}
\end{classroomqa}

\section{What a tachyon really means}

The name ``tachyon'' invites a misleading first thought. Start from the
relativistic dispersion relation with \(c=1\),

\begin{equation}
  E(p)=\sqrt{p^2+m^2}.
  \label{eq:l03-relativistic-dispersion}
\end{equation}

For a wave packet, the group velocity is

\begin{equation}
  v_g=\frac{dE}{dp}.
  \label{eq:l03-group-velocity}
\end{equation}

This reproduces familiar cases:

\begin{align}
  E_{\rm nr}(p)&=\frac{p^2}{2m},
  &v_g&=\frac{p}{m},
  \label{eq:l03-nr-group-velocity}\\
  E_\gamma(p)&=|p|,
  &|v_g|&=1,
  \label{eq:l03-photon-group-velocity}\\
  E_m(p)&=\sqrt{p^2+m^2},
  &|v_g|&=\frac{|p|}{\sqrt{p^2+m^2}}<1.
  \label{eq:l03-massive-group-velocity}
\end{align}

\begin{figure}[H]
  \centering
  \begin{minipage}[t]{0.48\textwidth}
    \centering
    \includegraphics[width=\textwidth]{lecture_03_figure_12.png}
    \par\smallskip\footnotesize (a) Nonrelativistic group velocity.
  \end{minipage}\hfill
  \begin{minipage}[t]{0.48\textwidth}
    \centering
    \includegraphics[width=\textwidth]{lecture_03_figure_13.png}
    \par\smallskip\footnotesize (b) Relativistic dispersion and its
    derivative.
  \end{minipage}
  \caption{Reading propagation speed from the slope of \(E(p)\).}
  \lectureframe{01:01:50 and 01:03:15}
  \label{fig:l03-group-velocity-board}
\end{figure}

If \(m^2=-\mu^2\), the formal expression becomes

\begin{equation}
  E(p)=\sqrt{p^2-\mu^2},
  \qquad
  v_g=\frac{p}{\sqrt{p^2-\mu^2}},
  \label{eq:l03-formal-tachyon-dispersion}
\end{equation}

which has \(|v_g|>1\) whenever the square root is real. But this is not the
right physical interpretation of the bosonic-string ground state. The
trouble appears already at low momentum, where the frequency is imaginary.

\subsection{Negative curvature of the field potential}

Consider a real scalar field in one spatial dimension,

\begin{equation}
  \left(\partial_t^2-\partial_x^2+m^2\right)\phi=0.
  \label{eq:l03-klein-gordon}
\end{equation}

For a plane wave \(\phi\propto e^{-i\omega t+ikx}\),

\begin{equation}
  \omega^2=k^2+m^2.
  \label{eq:l03-kg-dispersion}
\end{equation}

\begin{figure}[H]
  \centering
  \includegraphics[width=0.82\textwidth]{lecture_03_figure_14.png}
  \caption{The Klein--Gordon equation used to expose the instability behind
  negative mass squared.}
  \lectureframe{01:06:40}
  \label{fig:l03-klein-gordon-board}
\end{figure}

The conserved quadratic energy is

\begin{equation}
  \mathcal E
  =\frac12\int dx\,
  \left[
    \dot\phi^{\,2}
    +(\partial_x\phi)^2
    +m^2\phi^2
  \right].
  \label{eq:l03-field-energy}
\end{equation}

\begin{figure}[H]
  \centering
  \includegraphics[width=0.82\textwidth]{lecture_03_figure_15.png}
  \caption{The scalar-field energy, with the mass term supplying the local
  curvature of the potential.}
  \lectureframe{01:08:10}
  \label{fig:l03-field-energy-board}
\end{figure}

For \(m^2>0\), the local potential
\(V(\phi)=\tfrac12m^2\phi^2\) is a bowl. The configuration \(\phi=0\) sits at
its bottom. For \(m^2=-\mu^2\), the same point is the top of an inverted bowl.
Equation \eqref{eq:l03-kg-dispersion} becomes

\begin{equation}
  \omega^2=k^2-\mu^2.
  \label{eq:l03-tachyonic-modes}
\end{equation}

Modes with \(k^2<\mu^2\) have

\begin{equation}
  \omega=\pm i\sqrt{\mu^2-k^2},
  \qquad
  \phi_k(t)\propto
  e^{\pm\sqrt{\mu^2-k^2}\,t}.
  \label{eq:l03-growing-mode}
\end{equation}

One branch grows exponentially. The alleged vacuum is therefore unstable:
an arbitrarily small long-wavelength disturbance drives the field away from
\(\phi=0\). That is what a tachyon in the spectrum means.

\subsection{The upright-pendulum chain}

A row of coupled pendulums makes the distinction visible. Put every pendulum
upright. The uniform configuration exists, but it is balanced at a maximum
of potential energy. Nudge one pendulum and it falls, pulling on its
neighbors. A front of falling pendulums spreads through the chain at the
finite signal speed fixed by the couplings.

\begin{figure}[H]
  \centering
  \begin{minipage}[t]{0.48\textwidth}
    \centering
    \includegraphics[width=\textwidth]{lecture_03_figure_16.png}
    \par\smallskip\footnotesize (a) The coupled upright-pendulum model on the
    board.
  \end{minipage}\hfill
  \begin{minipage}[t]{0.48\textwidth}
    \centering
    \begin{tikzpicture}[>=Latex,scale=0.78]
      \draw[thick] (-0.3,0) -- (6.3,0);
      \foreach \x/\a in {0/0,1/0,2/8,3/28,4/55,5/72,6/78}{
        \fill (\x,0) circle (0.06);
        \draw[thick] (\x,0) -- ++({90-\a}:1.25);
        \fill[blue!65!black]
          ($(\x,0)+({90-\a}:1.25)$) circle (0.11);
      }
      \draw[red!70!black,->,thick] (1.7,-0.45) -- (4.6,-0.45)
        node[midway,below] {falling front};
    \end{tikzpicture}
  \end{minipage}
  \caption{Instability spreads causally through a system even though its
  linearized spectrum contains negative mass squared.}
  \lectureframe{01:10:40}
  \label{fig:l03-pendulum-instability}
\end{figure}

Nothing in this mechanism permits a controlled signal to outrun the causal
front. A tachyonic mode diagnoses a bad expansion point, not a stable
superluminal messenger.\editorialnote{For a local relativistic field
equation the front velocity remains the causal speed. The formal group
velocity of the real high-\(k\) branch does not turn the unstable vacuum into
a faster-than-light communication channel.}

\begin{classroomqa}
\audiencequestion{Is the massless state just found the graviton predicted by
string theory?\lecturetimestamp{01:14:27}}
\lecturerresponse{No. It is the two-polarization spin-one state of an open
string, so it is gauge-boson-like. The graviton is a spin-two excitation of a
closed string.}
\end{classroomqa}

\begin{classroomqa}
\audiencequestion{Is an upside-down potential mathematically allowed, and
where would the system finish after it falls?\lecturetimestamp{01:15:17}}
\lecturerresponse{It is a perfectly legitimate local approximation, but by
itself it has no ground state because the energy is unbounded below. A more
complete potential may bend upward again and possess a lower minimum. The
system then rolls or condenses toward that stable configuration.}
\end{classroomqa}

The ordinary bosonic string has precisely this defect. Its higher oscillator
states remain part of the spectrum, but they are excitations of an unstable
background. Consistent superstring theories remove the open-string tachyon in
their stable sectors while retaining massless and massive string
excitations.

\begin{classroomqa}
\audiencequestion{Do the other properties of the open-string state match a
photon as well as its two polarizations do?\lecturetimestamp{01:17:07}}
\lecturerresponse{The state behaves as a massless gauge boson. Its detailed
charge assignments, radiation rules, and identity depend on the string
model, so one should call it photon-like rather than identify it immediately
with electromagnetism. Those interaction properties are developed later.}
\end{classroomqa}

\section{How strings interact}

The free spectrum is only half the theory. Strings also join and split. It is
useful first to remember that an excited free string can carry angular
momentum. A straight rotating string turns about its center of mass, and its
energy and angular momentum rise together along the familiar stringy
trajectory.

\begin{figure}[H]
  \centering
  \includegraphics[width=0.72\textwidth]{lecture_03_figure_17.png}
  \caption{A rotating open string sketched about its center-of-mass point.}
  \lectureframe{01:19:30}
  \label{fig:l03-rotating-string-board}
\end{figure}

\begin{classroomqa}
\audiencequestion{Can a free string rotate end over end about one fixed
endpoint?\lecturetimestamp{01:20:02}}
\lecturerresponse{Not without an anchor. A free isolated system rotates about
its center of mass. If an endpoint is attached to another object or to a
brane, the combined system can rotate about a different center, just as an
electron and nucleus orbit their common center of mass while the much heavier
nucleus appears nearly fixed.}
\end{classroomqa}

\subsection{The open-string vertex}

For open strings the elementary interaction occurs when two endpoints meet
at the same spacetime point. There is an amplitude \(g_s\) for them to join
and form one string. Reversing the history gives an amplitude for one string
to split into two. Schematically,

\begin{equation}
  \text{open}+\text{open}
  \quad\rightleftarrows\quad
  \text{open}.
  \label{eq:l03-open-joining}
\end{equation}

\begin{figure}[H]
  \centering
  \begin{minipage}[t]{0.48\textwidth}
    \centering
    \includegraphics[width=\textwidth]{lecture_03_figure_18.png}
    \par\smallskip\footnotesize (a) The joining history on the board.
  \end{minipage}\hfill
  \begin{minipage}[t]{0.48\textwidth}
    \centering
    \begin{tikzpicture}[>=Latex,scale=0.9]
      \draw[->] (0,-2.1) -- (0,2.2) node[above] {time};
      \draw[blue!70!black,thick] (-2,-1.6) .. controls (-1.3,-0.9)
        and (-0.7,-0.35) .. (0,0);
      \draw[blue!70!black,thick] (2,-1.6) .. controls (1.3,-0.9)
        and (0.7,-0.35) .. (0,0);
      \draw[orange!85!black,thick] (0,0) .. controls (0.25,0.7)
        and (-0.15,1.25) .. (0.15,1.9);
      \fill[red!70!black] (0,0) circle (0.09);
      \node[right] at (0.12,0) {\(g_s\)};
    \end{tikzpicture}
  \end{minipage}
  \caption{Endpoint joining is the string analogue of a local particle
  interaction vertex; time reversal gives splitting.}
  \lectureframe{01:24:05}
  \label{fig:l03-open-vertex}
\end{figure}

The coupling is an amplitude. For one specified elementary process the
corresponding probability begins at order \(|g_s|^2\), with phase-space and
normalization factors supplied by the full calculation. The lecture moves
between the colloquial words ``probability'' and ``coupling''; keeping
amplitude and probability distinct prevents a spurious square.

The interaction does not depend on whether the external string happens to be
vibrating, rotating, or in some other oscillator state. Those states change
the detailed matrix element, but the same local joining rule is used.

\subsection{An open string can close itself}

Suppose the two endpoints of one open string meet. The joining rule cannot
know that they belong to the same string. They can fuse, producing a closed
loop:

\begin{equation}
  \text{open}\longrightarrow\text{closed}.
  \label{eq:l03-self-closing}
\end{equation}

\begin{figure}[H]
  \centering
  \begin{minipage}[t]{0.48\textwidth}
    \centering
    \includegraphics[width=\textwidth]{lecture_03_figure_19.png}
    \par\smallskip\footnotesize (a) The open string brings its endpoints
    together.
  \end{minipage}\hfill
  \begin{minipage}[t]{0.48\textwidth}
    \centering
    \begin{tikzpicture}[>=Latex,scale=0.92]
      \draw[blue!70!black,thick]
        (-2,0) .. controls (-1.2,1.5) and (1.2,1.5) .. (2,0);
      \fill[red!70!black] (-2,0) circle (0.08);
      \fill[red!70!black] (2,0) circle (0.08);
      \draw[->,thick] (-0.45,-0.45) -- (0.45,-0.45);
      \draw[orange!85!black,thick] (0,-1.55) ellipse (1.15 and 0.58);
      \node at (0,-2.25) {closed string};
    \end{tikzpicture}
  \end{minipage}
  \caption{Interacting open strings inevitably generate closed strings.}
  \lectureframe{01:27:20}
  \label{fig:l03-open-becomes-closed}
\end{figure}

This is not an optional embellishment. Any theory of interacting open
strings must contain closed-string states. The reverse process lets a closed
string open into an open string when the allowed endpoints and boundary
conditions are present.

\begin{classroomqa}
\audiencequestion{When strings fuse or split, does their oscillator energy
change, and does the theory determine how?\lecturetimestamp{01:28:05}}
\lecturerresponse{Yes. Energy and momentum are conserved for the complete
process, while excitation can be redistributed among the outgoing strings.
The interaction vertex and the external wavefunctions determine the
transition amplitude.}
\end{classroomqa}

\subsection{Closed strings can interact without opening}

The implication is one-way. Interacting open strings require closed strings,
but a theory may begin with only closed strings. Two closed loops can touch,
reconnect, and become one closed loop; the reverse history splits one loop
into two. No open endpoint need appear.

\begin{figure}[H]
  \centering
  \begin{minipage}[t]{0.48\textwidth}
    \centering
    \includegraphics[width=\textwidth]{lecture_03_figure_20.png}
    \par\smallskip\footnotesize (a) Closed-string fusion on the board.
  \end{minipage}\hfill
  \begin{minipage}[t]{0.48\textwidth}
    \centering
    \begin{tikzpicture}[>=Latex,scale=0.88]
      \draw[blue!70!black,thick] (-1.35,0) circle (0.75);
      \draw[blue!70!black,thick] (1.35,0) circle (0.75);
      \draw[->,thick] (-0.35,-1.15) -- (0.35,-1.15);
      \draw[orange!85!black,thick] (0,-2.3) ellipse (1.35 and 0.68);
      \node at (0,1.2) {closed \(+\) closed};
      \node at (0,-3.25) {closed};
    \end{tikzpicture}
  \end{minipage}
  \caption{Closed strings possess their own joining and splitting
  interaction and need not pass through an open-string state.}
  \lectureframe{01:29:05}
  \label{fig:l03-closed-fusion}
\end{figure}

\begin{classroomqa}
\audiencequestion{If a theory starts with closed strings, must it also
contain open strings?\lecturetimestamp{01:28:35}}
\lecturerresponse{No. Closed strings can interact by reconnection while
remaining closed throughout. Open strings require additional boundary data,
such as allowed endpoints or branes.}
\end{classroomqa}

\section{Why gravity cannot be avoided}

The closed-string oscillator calculation parallels the open-string one, but
there are independent left-moving and right-moving excitations. Their lowest
nontrivial matched level contains a massless spin-two state. That state has
the helicities of a graviton. Consequently every consistent interacting
closed-string theory contains gravity.

\begin{classroomqa}
\audiencequestion{Why does the lowest closed-string excitation have to be a
graviton, just as the open-string level had to be photon-like?
\lecturetimestamp{01:30:59}}
\lecturerresponse{The answer is the same kind of state counting, now with
left- and right-moving oscillator vectors. Their tensor product contains a
massless symmetric traceless spin-two polarization. The explicit closed-
string spectrum will be worked out next; at this stage the point is the
inevitability of that state.}
\end{classroomqa}

Open and closed strings may also interact with one another. A closed loop can
meet an open string and be absorbed, leaving a larger or differently excited
open string. Conversely an open string can radiate a closed loop. At the
particle level this is the emission or absorption of a graviton by matter.

\begin{figure}[H]
  \centering
  \begin{minipage}[t]{0.48\textwidth}
    \centering
    \includegraphics[width=\textwidth]{lecture_03_figure_21.png}
    \par\smallskip\footnotesize (a) Local reconnection of string segments.
  \end{minipage}\hfill
  \begin{minipage}[t]{0.48\textwidth}
    \centering
    \includegraphics[width=\textwidth]{lecture_03_figure_22.png}
    \par\smallskip\footnotesize (b) A closed loop attaches to an open
    string.
  \end{minipage}
  \caption{The board pictures behind mixed open--closed interactions and the
  universal possibility of graviton exchange.}
  \lectureframe{01:33:10 and 01:34:35}
  \label{fig:l03-mixed-interactions-board}
\end{figure}

\begin{figure}[H]
  \centering
  \begin{tikzpicture}[>=Latex,scale=0.88]
    \draw[blue!70!black,thick]
      (-3,-0.5) .. controls (-1.7,-0.1) and (-0.8,0.2) .. (0,0);
    \draw[blue!70!black,thick]
      (0,0) .. controls (0.9,0.3) and (2.0,0.1) .. (3,0.55);
    \draw[orange!85!black,thick] (-0.8,1.55) ellipse (0.75 and 0.45);
    \draw[->,red!70!black,thick] (-0.45,1.15) -- (-0.05,0.35);
    \node[blue!70!black] at (-2.2,-0.95) {open string};
    \node[orange!85!black] at (-0.8,2.25) {closed string};
    \node at (1.75,-0.5) {excited open string};
  \end{tikzpicture}
  \caption{A clean spacetime-free schematic of closed-string absorption by
  an open string.}
  \label{fig:l03-mixed-interaction-redraw}
\end{figure}

The graviton can couple to every string state because every string carries
energy and momentum and the joining rules allow closed-string exchange.
This is the string-theory origin of gravity's universality.

\begin{classroomqa}
\audiencequestion{Can a closed string become an open string, or can it join
an open string without violating the endpoint rule?\lecturetimestamp{01:31:20}}
\lecturerresponse{In a theory containing both sectors, a closed string can
attach to an open string and change its state. Whether a lone closed string
may simply open depends on whether the theory supplies permissible
endpoints. The local reconnection rule remains the same.}
\end{classroomqa}

\begin{classroomqa}
\audiencequestion{Can closed loops become permanently linked or
knotted?\lecturetimestamp{01:35:10}}
\lecturerresponse{At weak coupling strings are mostly transparent: they can
pass through one another unless the small reconnection amplitude acts. A
link or knot can therefore persist for some time, but topology can change
when an allowed string interaction occurs.}
\end{classroomqa}

\begin{classroomqa}
\audiencequestion{Is the interaction probability independent of the state
or speed of the string?\lecturetimestamp{01:36:09}}
\lecturerresponse{The local coupling is the same parameter, but the total
transition probability also depends on the wavefunctions and on how long
the relevant pieces remain in contact. A fast encounter generally offers
less interaction time than a slow one.}
\end{classroomqa}

\section{The scale and coupling still missing}

Two pieces of bookkeeping have deliberately been postponed. First, \(X\) is
a length while \(a_n\) is dimensionless, so the proportionality in
\eqref{eq:l03-mode-coordinates} must contain a length scale. Second, raising
the oscillator level cannot literally add a dimensionless number to a
physical mass squared. Restoring the string length \(\ell_s\), or equivalently
the Regge slope \(\alpha'\sim\ell_s^2\), supplies the missing dimensions:

\begin{equation}
  M^2\sim\frac{1}{\alpha'}(N-a).
  \label{eq:l03-dimensional-spectrum}
\end{equation}

The precise convention and normalization will be fixed in the next lecture.

\begin{classroomqa}
\audiencequestion{What fixes the unit of string length and the mass added by
one creation operator?\lecturetimestamp{01:37:03}}
\lecturerresponse{A dimensionful constant has been suppressed. It is the
string scale, usually written \(\ell_s\) or \(\alpha'\). The oscillator
operators are dimensionless; powers of that scale convert their level number
into physical length, tension, and mass squared.}
\end{classroomqa}

The second parameter is the string coupling \(g_s\). In the simplest
discussion it is one universal number controlling joining and splitting. In
richer models, endpoints can carry different labels and strings can meet at
junctions. A Y-shaped configuration is the string analogue of three flux
tubes meeting; in the old hadronic picture it is associated with a
three-quark baryon, whereas a single open segment resembles a quark--antiquark
meson.

\begin{classroomqa}
\audiencequestion{Can different string endpoints have different joining
probabilities?\lecturetimestamp{01:39:31}}
\lecturerresponse{The simplest model uses one coupling. More elaborate
theories attach charges or boundary labels to endpoints, and selection rules
then decide which joins are allowed. The detailed construction is deferred,
but the universal local joining idea survives.}
\end{classroomqa}

\begin{classroomqa}
\audiencequestion{Can a string split lengthwise or form a Y-shaped
junction?\lecturetimestamp{01:40:01}}
\lecturerresponse{Some string systems do admit junctions. They are not the
ordinary two-endpoint splitting process just drawn, but they occur when the
theory has compatible string charges. The meson--baryon flux-tube analogy is
a useful picture of the distinction.}
\end{classroomqa}

\begin{classroomqa}
\audiencequestion{Does changing the number of spacetime dimensions change
the interaction coupling?\lecturetimestamp{01:42:25}}
\lecturerresponse{Dimensions change the spectrum, phase space, and effective
couplings obtained after compactification. The fundamental string coupling
still organizes the perturbative expansion, but its lower-dimensional
appearance depends on the geometry and volume of the extra dimensions.}
\end{classroomqa}

\begin{classroomqa}
\audiencequestion{Is gravity what makes a free string move in the first
place?\lecturetimestamp{01:43:17}}
\lecturerresponse{No force is required to maintain motion. Initial position
and velocity specify a free solution; inertia carries it forward. Gravity or
another force changes that motion. Asking what caused the initial conditions
is a separate question, not part of the local equation of motion.}
\end{classroomqa}

We have therefore reached three linked conclusions. Quantizing an open
string produces a massless vector at its first excited level; the bosonic
ground state below it is unstable; and interactions inevitably introduce
closed strings, whose spectrum contains gravity. The next task is to restore
the missing scale and perform the closed-string quantization explicitly.
